\ifdefined\gkver\else\def\gkver{student}\fi
\documentclass[\gkver]{gaokaozhenti}
\newcommand{\ccnum}[1]{{\fontspec{Noto Serif CJK SC}#1}}
\begin{document}

\chapter{2011年全国新课标}
%% paperType: 理综物理部分

\begin{enumerate}
\item
%% number: 14
%% paperName: 2011年全国新课标第 14 题 6 分
%% typeId: 1
%% type: 单选题
%% chapter: 磁场
%% point: 电流的磁场
%% method: 等效替代
%% score: 6
%% degree: 650
%% duplicateId: 0
%% body:
为了解释地球的磁性，19 世纪安培假设：地球的磁场是由绕过地心的轴的环形电流 $I$ 引起的。在下列四个图中，正确表示安培假设中环形电流方向的是\xzanswer{B}

\fourchoices[answer=B, ispicture=true, h=2.6cm]
{figs/14a.png}
{figs/14b.png}
{figs/14c.png}
{figs/14d.png}

%% answer: B
%% memo:
\memoanswer{
	根据地磁场分布和安培定则判断可知正确答案是 B。
}

\item
%% number: 15
%% paperName: 2011年全国新课标第 15 题 6 分
%% typeId: 2
%% type: 多选题
%% chapter: 功和能
%% point: 动能定理
%% method: 无
%% score: 6
%% degree: 450
%% duplicateId: 0
%% body:
一质点开始时做匀速直线运动，从某时刻起受到一恒力作用。此后，该质点的动能可能\xzanswer{ABD}

\fourchoices[answer=ABD]
{一直增大}
{先逐渐减小至零，再逐渐增大}
{先逐渐增大至某一最大值，再逐渐减小}
{先逐渐减小至某一非零的最小值，再逐渐增大}

%% answer: ABD
%% memo:
\memoanswer{
	当恒力方向与速度在一条直线上，质点的动能可能一直增大，也可能先逐渐减小至零，再逐渐增大。当恒力方向与速度不在一条直线上，质点的动能可能一直增大，也可能先逐渐减小至某一非零的最小值，再逐渐增大。所以正确答案是 ABD。
}

\item
%% number: 16
%% paperName: 2011年全国新课标第 16 题 6 分
%% typeId: 2
%% type: 多选题
%% chapter: 功和能
%% point: 机械能守恒定律
%% method: 无
%% score: 6
%% degree: 450
%% duplicateId: 0
%% body:
一蹦极运动员身系弹性蹦极绳从水面上方的高台下落，到最低点时距水面还有数米距离。假定空气阻力可忽略，运动员可视为质点，下列说法正确的是\xzanswer{ABC}

\fourchoices[answer=ABC]
{运动员到达最低点前重力势能始终减小}
{蹦极绳张紧后的下落过程中，弹性力做负功，弹性势能增加}
{蹦极过程中，运动员、地球和蹦极绳所组成的系统机械能守恒}
{蹦极过程中，重力势能的改变与重力势能零点的选取有关}

%% answer: ABC
%% memo:
\memoanswer{
	运动员到达最低点过程中，重力做正功，所以重力势能始终减少，A 项正确。蹦极绳张紧后的下落过程中，弹性力做负功，弹性势能增加，B 项正确。蹦极过程中，运动员、地球和蹦极绳所组成的系统，只有重力和弹性力做功，所以机械能守恒，C 项正确。重力势能的改变与重力势能零点选取无关，D 项错误。
}

\item
%% number: 17
%% paperName: 2011年全国新课标第 17 题 6 分
%% typeId: 1
%% type: 单选题
%% chapter: 交流电
%% point: 变压器
%% method: 无
%% score: 6
%% degree: 650
%% duplicateId: 0
%% body:
如图，一理想变压器原副线圈的匝数比为 $1:2$；副线圈电路中接有灯泡，灯泡的额定电压为 $220\UV$，额定功率为 $22\UW$；原线圈电路中接有电压表和电流表。现闭合开关，灯泡正常发光。若用 $U$ 和 $I$ 分别表示此时电压表和电流表的读数，则\xzanswer{A}

\onepicture[w=3.4cm]{figs/17.png}

\fourchoices[answer=A]
{$U=110\UV$，$I=0.2\UA$}
{$U=110\UV$，$I=0.05\UA$}
{$U=110\sqrt{2}\UV$，$I=0.2\UA$}
{$U=110\sqrt{2}\UV$，$I=0.2\sqrt{2}\UA$}

%% answer: A
%% memo:
\memoanswer{
	$U_{2}=220\UV$，根据 $U_{1}:U_{2}=n_{1}:n_{2}$ 得，$U_{1}=110\UV$。$I_{2}=P/U_{2}=0.1\UA$，根据 $I_{1}:I_{2}=n_{2}:n_{1}$ 得 $I_{1}=0.2\UA$。所以正确答案是 A。
}

\item
%% number: 18
%% paperName: 2011年全国新课标第 18 题 6 分
%% typeId: 2
%% type: 多选题
%% chapter: 磁场
%% point: 安培力
%% method: 无
%% score: 6
%% degree: 450
%% duplicateId: 0
%% body:
电磁轨道炮工作原理如图所示。待发射弹体可在两平行轨道之间自由移动，并与轨道保持良好接触。电流 $I$ 从一条轨道流入，通过导电弹体后从另一条轨道流回。轨道电流可形成在弹体处垂直于轨道面的磁场（可视为匀强磁场），磁感应强度的大小与 $I$ 成正比。通电的弹体在轨道上受到安培力的作用而高速射出。现欲使弹体的出射速度增加至原来的 2 倍，理论上可采用的办法是\xzanswer{BD}

\onepicture[w=4cm]{figs/18.png}

\fourchoices[answer=BD]
{只将轨道长度 $L$ 变为原来的 2 倍}
{只将电流 $I$ 增加至原来的 2 倍}
{只将弹体质量减至原来的一半}
{将弹体质量减至原来的一半，轨道长度 $L$ 变为原来的 2 倍，其它量不变}

%% answer: BD
%% memo:
\memoanswer{
	利用动能定理有 $BIlL=\frac{1}{2}mv^{2}$，$B=kI$，解得 $v=\sqrt{\frac{2kI^{2}lL}{m}}$（$l$ 为两轨道间距）。所以正确答案是 BD。
}

\item
%% number: 19
%% paperName: 2011年全国新课标第 19 题 6 分
%% typeId: 1
%% type: 单选题
%% chapter: 曲线运动
%% point: 开普勒定律
%% method: 无
%% score: 6
%% degree: 450
%% duplicateId: 0
%% body:
卫星电话信号需要通地球同步卫星传送。如果你与同学在地面上用卫星电话通话，则从你发出信号至对方接收到信号所需最短时间最接近于（可能用到的数据：月球绕地球运动的轨道半径约为 $3.8\times10^{5}\Ukm$，运行周期约为 27 天，地球半径约为 $6400\Ukm$，无线电信号的传播速度为 $3\times10^{8}\Ums$。）\xzanswer{B}

\fourchoices[answer=B]
{$0.1\Us$}
{$0.25\Us$}
{$0.5\Us$}
{$1\Us$}

%% answer: B
%% memo:
\memoanswer{
	同步卫星和月球都是地球的卫星，$r^{3}\propto T^{2}$，因此同步卫星的轨道半径是地月距离的 $\frac{1}{9}$ 约为 $42000\Ukm$，同步卫星离地面高度约为 $36000\Ukm$，电磁波往返一次经历时间约为 $(3.6\times10^{7}\times2)\div(3\times10^{8})\Us=0.24\Us$。
}

\item
%% number: 20
%% paperName: 2011年全国新课标第 20 题 6 分
%% typeId: 1
%% type: 单选题
%% chapter: 电场
%% point: 带电粒子的偏转
%% method: 无
%% score: 6
%% degree: 450
%% duplicateId: 0
%% body:
一带负电荷的质点，在电场力作用下沿曲线 abc 从 a 运动到 c，已知质点的速率是递减的。关于 b 点电场强度 $E$ 的方向，下列图示中可能正确的是（虚线是曲线在 b 点的切线）\xzanswer{D}

\fourchoices[answer=D, ispicture=true, h=2.4cm]
{figs/20a.png}
{figs/20b.png}
{figs/20c.png}
{figs/20d.png}

%% answer: D
%% memo:
\memoanswer{
	主要考查电场力方向和曲线运动所受合外力与轨迹的关系。正确答案是 D。
}

\item
%% number: 21
%% paperName: 2011年全国新课标第 21 题 6 分
%% typeId: 1
%% type: 单选题
%% chapter: 运动和力
%% point: 传送带
%% method: 无
%% score: 6
%% degree: 450
%% duplicateId: 0
%% body:
如图，在光滑水平面上有一质量为 $m_{1}$ 的足够长的木板，其上叠放一质量为 $m_{2}$ 的木块。假定木块和木板之间的最大静摩擦力和滑动摩擦力相等。现给木块施加一随时间 $t$ 增大的水平力 $F=kt$（$k$ 是常数），木板和木块加速度的大小分别为 $a_{1}$ 和 $a_{2}$，下列反映 $a_{1}$ 和 $a_{2}$ 变化的图线中正确的是\xzanswer{A}

\onepicture[w=4.5cm]{figs/21a.png}

\fourchoices[answer=A, ispicture=true, h=2.2cm]
{figs/21b.png}
{figs/21c.png}
{figs/21d.png}
{figs/21e.png}

%% answer: A
%% memo:
\memoanswer{
	木块和木板之间相对静止时，所受的摩擦力为静摩擦力。在达到最大静摩擦力前，木块和木板以相同加速度运动，根据牛顿第二定律 $a_{1}=a_{2}=\frac{kt}{m_{1}+m_{2}}$，木块和木板相对运动时，$a_{1}=\frac{\mu m_{2}g}{m_{1}}$ 恒定不变，$a_{2}=\frac{kt}{m_{2}}-\mu g$。所以正确答案是 A。
}

\item
%% number: 22
%% paperName: 2011年全国新课标第 22 题 5 分
%% typeId: 4
%% type: 实验题
%% chapter: 电路
%% point: 电路实验
%% method: 等效替代
%% score: 5
%% degree: 650
%% duplicateId: 0
%% body:
为了测量一微安表头 A 的内阻，某同学设计了如图所示的电路。图中，A$_{0}$ 是标准电流表，$R_{0}$ 和 $R_{N}$ 分别是滑动变阻器和电阻箱，S 和 S$_{1}$ 分别是单刀双掷开关和单刀开关，$E$ 是电池。完成下列实验步骤中的填空：
\begin{enumerate}
	\item
	将 S 拨向接点 1，接通 S$_{1}$，调节\tkanswer[1.4cm]{$R_{0}$}，使待测表头指针偏转到适当位置，记下此时\tkanswer[3cm]{标准电流表 A$_{0}$}的读数 $I$；
	\item
	然后将 S 拨向接点 2，调节\tkanswer[1.4cm]{$R_{N}$}，使\tkanswer[3.6cm]{标准电流表 A$_{0}$ 的示数为 $I$}，记下此时 $R_{N}$ 的读数；
	\item
	多次重复上述过程，计算 $R_{N}$ 读数的\tkanswer[1.6cm]{平均值}，此即为待测微安表头内阻的测量值。
\end{enumerate}

\onepicture[w=6.4cm, align=right]{figs/22.png}

%% answer: 
\jdanswer{
\begin{enumerate}
	\item
	$R_{0}$、标准电流表 A$_{0}$；
	\item
	$R_{N}$，使标准电流表 A$_{0}$ 的示数为 $I$；
	\item
	平均值。
\end{enumerate}
}
%% memo:
\memoanswer{
	本题原卷及材料中未提供详解，待补充。
}

\item
%% number: 23
%% paperName: 2011年全国新课标第 23 题 10 分
%% typeId: 4
%% type: 实验题
%% chapter: 直线运动
%% point: DIS测位移平均速度加速度
%% method: 无
%% score: 10
%% degree: 450
%% duplicateId: 0
%% body:
利用图 \ref{2011全国新课标23a} 所示的装置可测量滑块在斜面上运动的加速度。一斜面上安装有两个光电门，其中光电门乙固定在斜面上靠近底端处，光电门甲的位置可移动，当一带有遮光片的滑块自斜面上滑下时，与两个光电门都相连的计时器可以显示出遮光片从光电门甲至乙所用的时间 $t$。改变光电门甲的位置进行多次测量，每次都使滑块从同一点由静止开始下滑，并用米尺测量甲、乙之间的距离 $s$，记下相应的 $t$ 值；所得数据如下表所示。

\begin{center}\small
\begin{tblr}{|c|c|c|c|c|c|c|}
\hline
$s(\Um)$ & 0.500 & 0.600 & 0.700 & 0.800 & 0.900 & 0.950 \\
\hline
$t(\UmsT)$ & 292.9 & 371.5 & 452.3 & 552.8 & 673.8 & 776.4 \\
\hline
$s/t(\Ums)$ & 1.71 & 1.62 & 1.55 & 1.45 & 1.34 & 1.22 \\
\hline
\end{tblr}
\end{center}

完成下列填空和作图：
\begin{enumerate}
	\item
	若滑块所受摩擦力为一常量，滑块加速度的大小 $a$、滑块经过光电门乙时的瞬时速度 $v_{1}$、测量值 $s$ 和 $t$ 四个物理量之间所满足的关系式是\tkanswer[3.2cm]{$s=v_{1}t-\frac{1}{2}at^{2}$}；
	\item
	根据表中给出的数据，在图 \ref{2011全国新课标23b} 给出的坐标纸上画出 $\frac{s}{t}-t$ 图线；
	\item
	由所画出的 $\frac{s}{t}-t$ 图线，得出滑块加速度的大小为 $a=$ \tkanswer[1.6cm]{$2.0$} $\Umsq$（保留 2 位有效数字）。
\end{enumerate}

\twopicture[labelA=2011全国新课标23a, labelB=2011全国新课标23b, align=right, wA=5cm, wB=5cm]
{figs/23a.png}
{figs/23b.png}

%% answer: 
\jdanswer{
\begin{enumerate}
	\item
	$s=v_{1}t-\frac{1}{2}at^{2}$；
	\item
	见下图；
	\item
	$2.0\Umsq$。
\end{enumerate}
}
%% memo:
\memoanswer{
	（1）由 $s=v_{0}t+\frac{1}{2}at^{2}$，$v_{1}=v_{0}+at$ 得 $s=v_{1}t-\frac{1}{2}at^{2}$。

	（2）见下图。

	\onepicture[w=5cm]{figs/23_an.png}

	（3）作图求出斜率 $k=-0.9897\Umsq$，$a=2|k|\approx2.0\Umsq$。
}

\item
%% number: 24
%% paperName: 2011年全国新课标第 24 题 13 分
%% typeId: 6
%% type: 计算题
%% chapter: 直线运动
%% point: 匀变速直线运动
%% method: 无
%% score: 13
%% degree: 650
%% duplicateId: 0
%% body:
甲乙两辆汽车都从静止出发做加速直线运动，加速度方向一直不变。在第一段时间间隔内，两辆汽车的加速度大小不变，汽车乙的加速度大小是甲的两倍；在接下来的相同时间间隔内，汽车甲的加速度大小增加为原来的两倍，汽车乙的加速度大小减小为原来的一半。求甲乙两车各自在这两段时间间隔内走过的总路程之比。

%% answer: 
\jdanswer{
	$\frac{5}{7}$
}
%% memo:
\memoanswer{
	设甲开始的加速度为 $a$，两段时间间隔都为 $t$，则甲在两段时间内的总路程为：

	$s_{\text{甲}}=\frac{1}{2}at^{2}+(at)t+\frac{1}{2}(2a)t^{2}=2.5at^{2}$，

	乙在两段时间内的总路程为：

	$s_{\text{乙}}=\frac{1}{2}(2a)t^{2}+(2at)t+\frac{1}{2}at^{2}=3.5at^{2}$，

	由上两式得：$\frac{s_{\text{甲}}}{s_{\text{乙}}}=\frac{5}{7}$。
}

\item
%% number: 25
%% paperName: 2011年全国新课标第 25 题 19 分
%% typeId: 6
%% type: 计算题
%% chapter: 磁场
%% point: 带电粒子在磁场中的运动
%% method: 无
%% score: 19
%% degree: 450
%% duplicateId: 0
%% body:
如图，在区域Ⅰ（$0\leq x\leq d$）和区域Ⅱ（$d<x\leq2d$）内分别存在匀强磁场，磁感应强度大小分别为 $B$ 和 $2B$，方向相反，且都垂直于 $Oxy$ 平面。一质量为 $m$、带电荷量 $q$（$q>0$）的粒子 a 于某时刻从 $y$ 轴上的 $P$ 点射入区域Ⅰ，其速度方向沿 $x$ 轴正向。已知 a 在离开区域Ⅰ时，速度方向与 $x$ 轴正方向的夹角为 $30^{\circ}$；此时，另一质量和电荷量均与 a 相同的粒子 b 也从 $P$ 点沿 $x$ 轴正向射入区域Ⅰ，其速度大小是 a 的 $\frac{1}{3}$。不计重力和两粒子之间的相互作用力。求：
\begin{enumerate}
	\item
	粒子 a 射入区域Ⅰ时速度的大小；
	\item
	当 a 离开区域Ⅱ时，a、b 两粒子的 $y$ 坐标之差。
\end{enumerate}

\onepicture[w=4.5cm, align=right]{figs/25.png}

%% answer: 
\jdanswer{
\begin{enumerate}
	\item
	$\frac{2qBd}{m}$
	\item
	$\frac{2}{3}(\sqrt{3}-2)d$
\end{enumerate}
}
%% memo:
\memoanswer{
	（1）设粒子 a 在Ⅰ内做匀速圆周运动的圆心为 C（在 $y$ 轴上），半径为 $R_{a1}$，粒子速率为 $v_{a}$，运动轨迹与两磁场区域边界的交点为 $P'$，如图，由洛伦兹力公式和牛顿第二定律得

	$qv_{a1}B=m\frac{v_{a1}^{2}}{R_{a1}}$\ccnum{①}

	由几何关系得 $\angle POP'=\theta$\ccnum{②}

	$R_{a1}=\frac{d}{\sin\theta}$\ccnum{③}

	式中，$\theta=30^{\circ}$，由\ccnum{①}\ccnum{②}\ccnum{③}式得 $v_{a1}=\frac{2qBd}{m}$\ccnum{④}

	（2）设粒子 a 在Ⅱ内做圆周运动的圆心为 $O_{a}$，半径为 $R_{a2}$，射出点为 $P_{a}$（图中未画出轨迹），$\angle P'O_{a}P_{a}=\theta'$。由洛伦兹力公式和牛顿第二定律得

	$qv_{a}(2B)=m\frac{v_{a}^{2}}{R_{a2}}$\ccnum{⑤}

	由\ccnum{①}\ccnum{⑤}式得 $R_{a2}=\frac{R_{a1}}{2}$\ccnum{⑥}

	C、$P'$ 和 $O_{a}$ 三点共线，且由\ccnum{⑥}式知 $O_{a}$ 点必位于 $x=\frac{3}{2}d$\ccnum{⑦} 的平面上。由对称性知，$P_{a}$ 点与 $P'$ 点纵坐标相同，即 $y_{P_{a}}=R_{a1}\cos\theta+h$\ccnum{⑧}，式中，$h$ 是 C 点的 $y$ 坐标。

	\onepicture[w=6.5cm]{figs/25_an.png}

	设 b 在Ⅰ中运动的轨道半径为 $R_{b1}$，由洛伦兹力公式和牛顿第二定律得 $q\left(\frac{v_{a}}{3}\right)B=m\frac{\left(\frac{v_{a}}{3}\right)^{2}}{R_{b1}}$\ccnum{⑨}

	设 a 到达 $P_{a}$ 点时，b 位于 $P_{b}$ 点，转过的角度为 $\alpha$。如果 b 没有飞出Ⅰ，则

	$\frac{t}{T_{a2}}=\frac{\theta'}{2\pi}$\ccnum{⑩}，$\frac{t}{T_{b1}}=\frac{\alpha}{2\pi}$\ccnum{⑪}

	式中，$t$ 是 a 在区域Ⅱ中运动的时间，而

	$T_{a2}=\frac{2\pi R_{a2}}{v}$\ccnum{⑫}，$T_{b1}=\frac{2\pi R_{b1}}{v/3}$\ccnum{⑬}

	由\ccnum{⑤}\ccnum{⑨}\ccnum{⑩}\ccnum{⑪}\ccnum{⑫}\ccnum{⑬}式得 $\alpha=30^{\circ}$\ccnum{⑭}

	由\ccnum{①}\ccnum{③}\ccnum{⑨}\ccnum{⑩}式可见，b 没有飞出。$P_{b}$ 点的 $y$ 坐标为 $y_{P_{b}}=R_{b1}(2+\cos\theta)+h$\ccnum{⑮}

	由\ccnum{①}\ccnum{③}\ccnum{⑧}\ccnum{⑨}\ccnum{⑩}式及题给条件得，a、b 两粒子的 $y$ 坐标之差为

	$y_{P_{a}}-y_{P_{b}}=\frac{2}{3}(\sqrt{3}-2)d$\ccnum{⑯}
}

\item
%% number: 33
%% paperName: 2011年全国新课标第 33 题 9 分
%% typeId: 6
%% type: 计算题
%% chapter: 气体性质
%% point: 玻意耳定律
%% method: 无
%% score: 9
%% degree: 450
%% duplicateId: 0
%% body:
【物理——选修 3-3】（15 分）

（1）（6 分）对于一定质量的理想气体，下列说法正确的是\xzanswer{ADE}。

\fivechoices[answer=ADE]
{若气体的压强和体积都不变，其内能也一定不变}
{若气体的内能不变，其状态也一定不变}
{若气体的温度随时间不断升高，其压强也一定不断增大}
{气体温度每升高 $1\UK$ 所吸收的热量与气体经历的过程有关}
{当气体温度升高时，气体的内能一定增大}

（2）（9 分）如图，一上端开口、下端封闭的细长玻璃管，下部有长 $l_{1}=66\Ucm$ 的水银柱，中间封有长 $l_{2}=6.6\Ucm$ 的空气柱，上部有长 $l_{3}=44\Ucm$ 的水银柱，此时水银面恰好与管口平齐。已知大气压强为 $p_{0}=76\UcmHg$。如果使玻璃管绕底端在竖直平面内缓慢地转动一周，求在开口向下和转回到原来位置时管中空气柱的长度。封入的气体可视为理想气体，在转动过程中没有发生漏气。

\onepicture[h=6cm, align=right]{figs/33.png}

%% answer: 
\jdanswer{
	（2）$12\Ucm$，$9.2\Ucm$。
}
%% memo:
\memoanswer{
	（1）理想气体的内能只由温度决定，由理想气体状态方程 $\frac{PV}{T}=C$ 可知，若气体的压强和体积都不变，温度 $T$ 也不变，所以内能也一定不变，A、E 选项正确。若气体的内能不变，则温度 $T$ 不变，但气体的压强和体积可以改变，B 项错误。若气体的温度升高，体积增大，其压强可以不变，C 项错误。由热力学第一定律 $\Delta U=Q+W$ 知，D 选项正确。

	（2）设玻璃管开口向上时，空气柱压强为 $P_{1}=P_{0}+\rho gl_{3}$\ccnum{①}

	式中，$\rho$ 和 $g$ 分别表示水银的密度和重力加速度。

	玻璃管开口向下时，原来上部的水银有一部分会流出，封闭端会有部分真空。设此时开口端剩下的水银柱长度为 $x$，则 $P_{2}=\rho gl_{1}$，$P_{2}+\rho gx=P_{0}$\ccnum{②}，式中，$P_{2}$ 为管内空气柱的压强。

	由玻意耳定律得 $P_{1}(Sl_{2})=P_{2}(Sh)$\ccnum{③}，式中，$h$ 是此时空气柱的长度，$S$ 为玻璃管的横截面积。

	由\ccnum{①}\ccnum{②}\ccnum{③}式和题给条件得 $h=12\Ucm$\ccnum{④}

	从开始转动一周后，设空气柱的压强为 $P_{3}$，则 $P_{3}=P_{0}+\rho gx$\ccnum{⑤}

	由玻意耳定律得 $P_{1}(Sl_{2})=P_{3}(Sh')$\ccnum{⑥}，式中，$h'$ 是此时空气柱的长度。

	由\ccnum{①}\ccnum{②}\ccnum{③}\ccnum{⑤}\ccnum{⑥}式得 $h'\approx9.2\Ucm$\ccnum{⑦}。
}

\item
%% number: 34
%% paperName: 2011年全国新课标第 34 题 9 分
%% typeId: 6
%% type: 计算题
%% chapter: 光学
%% point: 光的折射
%% method: 无
%% score: 9
%% degree: 450
%% duplicateId: 0
%% body:
【物理——选修 3-4】（15 分）

（1）（6 分）一振动周期为 $T$，振幅为 $A$，位于 $x=0$ 点的波源从平衡位置沿 $y$ 轴正向开始做简谐振动，该波源产生的一维简谐横波沿 $x$ 轴正向传播，波速为 $v$，传播过程中无能量损失，一段时间后，该振动传播至某质点 P，关于质点 P 振动的说法正确的是\xzanswer{ABE}。

\fivechoices[answer=ABE]
{振幅一定为 $A$}
{周期一定为 $T$}
{速度的最大值一定为 $v$}
{开始振动的方向沿 $y$ 轴向上或向下取决于它离波源的距离}
{若 P 点与波源距离 $s=vT$，则质点 P 的位移与波源的相同}

（2）（9 分）一半圆柱形透明物体横截面如图所示，底面 AOB 镀银（图中粗线），O 表示半圆截面的圆心。一束光线在横截面内从 M 点的入射角为 $30^{\circ}$，$\angle MOA=60^{\circ}$，$\angle NOB=30^{\circ}$。求：
\begin{enumerate}
	\item
	光线在 M 点的折射角；
	\item
	透明物体的折射率。
\end{enumerate}

\onepicture[w=5cm, align=right]{figs/34.png}

%% answer: 
\jdanswer{
\begin{enumerate}
	\item
	$15^{\circ}$
	\item
	$\frac{\sqrt{6}+\sqrt{2}}{2}$
\end{enumerate}
}
%% memo:
\memoanswer{
	（1）由波的形成与传播可知，正确答案是 ABE。

	（2）（i）解：如图，透明物体内部的光路为折线 MPN，Q、M 点相对于底面 EF 对称，Q、P 和 N 三点共线。

	设在 M 点处，光的入射角为 $i$，折射角为 $r$，$\angle OMQ=\alpha$，$\angle PNF=\beta$。

	根据题意有 $\alpha=30^{\circ}$\ccnum{①}

	由几何关系得，$\angle PNO=\angle PQO=r$，于是

	$\beta+r=30^{\circ}$\ccnum{②}

	且 $\alpha+r=\beta$\ccnum{③}

	由\ccnum{①}\ccnum{②}\ccnum{③}式得 $r=15^{\circ}$\ccnum{④}

	\onepicture[w=6cm]{figs/34_an.png}

	（ii）根据折射率公式有 $n=\frac{\sin i}{\sin\gamma}$\ccnum{⑤}

	由\ccnum{④}\ccnum{⑤}式得 $n=\frac{\sqrt{6}+\sqrt{2}}{2}$。
}

\item
%% number: 35
%% paperName: 2011年全国新课标第 35 题 9 分
%% typeId: 6
%% type: 计算题
%% chapter: 运动和力
%% point: 动量和能量
%% method: 无
%% score: 9
%% degree: 450
%% duplicateId: 0
%% body:
【物理——选修 3-5】（15 分）

（1）（6 分）在光电效应实验中，某金属的截止频率相应的波长为 $\lambda_{0}$，该金属的逸出功为\tkanswer[2.4cm]{$h\frac{c}{\lambda_{0}}$}。若用波长为 $\lambda$（$\lambda<\lambda_{0}$）的单色光做该实验，则其遏止电压为\tkanswer[3.4cm]{$\frac{hc}{e}\frac{\lambda_{0}-\lambda}{\lambda_{0}\lambda}$}。已知电子的电荷量、真空中的光速和普朗克常量分别为 $e$、$c$ 和 $h$。

（2）（9 分）如图，A、B、C 三个木块的质量均为 $m$。置于光滑的水平面上，B、C 之间有一轻质弹簧，弹簧的两端与木块接触可不固连。将弹簧压紧到不能再压缩时用细线把 B 和 C 紧连，使弹簧不能伸展，以至于 B、C 可视为一个整体。现 A 以初速 $v_{0}$ 沿 B、C 的连线方向朝 B 运动，与 B 相碰并粘合在一起。以后细线突然断开，弹簧伸展，从而使 C 与 A、B 分离。已知 C 离开弹簧后的速度恰为 $v_{0}$。求弹簧释放的势能。

\onepicture[w=7cm, align=right]{figs/35.png}

%% answer: 
\jdanswer{
	$\frac{1}{3}mv_{0}^{2}$
}
%% memo:
\memoanswer{
	（1）由 $w_{\text{逸}}=h\nu_{\text{极限}}$ 和 $\nu=\frac{c}{\lambda}$ 得 $w_{\text{逸}}=h\frac{c}{\lambda_{0}}$。

	由爱因斯坦质能方程 $E_{k}=h\frac{c}{\lambda}-w_{\text{逸}}$ 和 $E_{k}=Ue$ 得 $U_{\text{截止}}=\frac{hc}{e}\frac{\lambda_{0}-\lambda}{\lambda_{0}\lambda}$。

	（2）设碰后 A、B 和 C 的共同速度的大小为 $v$，由动量守恒得 $3mv=mv_{0}$\ccnum{①}

	设 C 离开弹簧时，A、B 的速度大小为 $v_{1}$，由动量守恒得 $3mv=2mv_{1}+mv_{0}$\ccnum{②}

	设弹簧的弹性势能为 $E_{p}$，从细线断开到 C 与弹簧分开的过程中机械能守恒，有

	$\frac{1}{2}(3m)v^{2}+E_{p}=\frac{1}{2}(2m)v_{1}^{2}+\frac{1}{2}mv_{0}^{2}$\ccnum{③}

	由\ccnum{①}\ccnum{②}\ccnum{③}式得弹簧所释放的势能为 $E_{p}=\frac{1}{3}mv_{0}^{2}$\ccnum{④}。
}

\end{enumerate}

\end{document}
